\chapter{Conclusions}
\label{ch:conclusions}

The Framingham case study demonstrates that the \gls{bra} implementation can expose operational tradeoffs that are difficult to see from the existing route plan alone. For the currently assigned bus riders, \gls{bra} serves the same 4,780 students using 48 buses rather than 51, a reduction of 3 buses, or approximately 5.9\%. At the per-bus daily lease cost discussed in \hyperref[ap:stakeholder]{Appendix~\ref*{ap:stakeholder}} (low hundreds of dollars), this corresponds to a meaningful annual lease-cost reduction before accounting for fuel and labor. At the same time, estimated mean travel time falls from 43.08 minutes to 30.52 minutes, and the maximum observed travel time falls from 133.38 minutes to 58.31 minutes. Under the modeled assumptions and conservative operational constraints, this suggests that the current routing pattern is improved upon by feasible alternatives that use fewer vehicles while also reducing student ride times---a more efficient and equitable system.

The all-students experiments show that expanding transportation access creates a different tradeoff. Serving students who live at least 1.5 miles from school is feasible under the modeled assumptions, though closer to the fleet limit. In \autoref{tab:travel_times_all}, the model uses 62 buses for the 1.5-mile eligibility case, leaving no students unassigned under the assumed constraints. When all students are considered, the model serves a larger total population but leaves 13 students unassigned, indicating that universal default service is not achievable without either additional resources, relaxed service standards, revised stop assignment rules, or a different policy definition of eligibility.

These results also illustrate the value of the Pareto-front framing. Rather than producing a single ``best'' routing plan, the co-design formulation can show the frontier between competing objectives. This is useful for planning because the preferred solution may depend on institutional priorities: one configuration may minimize buses, another maximizes student coverage, and another could reduce long rides for vulnerable groups. By placing the current \gls{fps} routing plan on the same frontier as generated alternatives, the framework allows planners to distinguish between tradeoffs that are structurally unavoidable and inefficiencies that can be improved through redesign.

Equity considerations are central to this interpretation. The results show that students from non-English-speaking households, non-car-owning households, and students with disabilities all experience lower estimated mean travel times under \gls{bra} than under the current routing baseline. This suggests that efficiency improvements do not come at the expense of vulnerable groups. However, the current experiments evaluate equity primarily after routing, rather than embedding equity constraints directly in the optimization objective. Future versions should incorporate explicit constraints or penalties for unequal service outcomes, such as subgroup-specific maximum ride times, minimum coverage rates, or limits on the spatial concentration of unassigned students. Additionally, inspired by the work of \textcite{vanweeEvaluatingTransportEquity2021}, promising outcomes could be gained by identifying ways of indexing inequity levels and incorporating them into \gls{bra}.

The results should therefore be interpreted as planning evidence rather than deployment-ready schedules. The current model assumes a constant travel speed, does not include intersection delay or time-varying congestion, and lacks complete student-level linkage between home addresses, current stops, bus assignments, disability status, and wheelchair-accessibility requirements. These limitations affect the precision of individual travel-time estimates and may change which route chains are feasible in practice. Before implementation, the model should be recalibrated using district-validated stop assignments, complete accessibility data, and travel-time estimates derived from \gls{avl} data or a time-dependent routing engine.

Overall, the case study supports the central claim of this work: a co-design-based routing framework can help districts reason about bus network redesign under multiple objectives. \Gls{bra} provides a practical backend for generating feasible routing alternatives, while the Pareto-front representation clarifies how operational efficiency, coverage expansion, and equity goals interact. For \gls{fps}, the preliminary results indicate that meaningful improvements may be possible, particularly for currently assigned riders, but that broader opt-out service requires explicit policy choices about eligibility thresholds, fleet deployment, maximum ride times, and equity guarantees.

Taken together, this thesis contributes a compositional framework (grounded in the monotone co-design theory of \textcite{censi2016Mathematical}) for analyzing school bus routing as a co-design problem rather than a standalone vehicle-routing task. The \gls{mdpi} decomposition (fleet, monitor, driver, fuel, maintenance, and routing service components composed via shared functionality/resource interfaces, as in \autoref{fig:codesign}) is what enables planners to vary one component at a time and analyze the resulting Pareto front, rather than re-solving an entire monolithic \gls{milp} for each scenario. By adapting the \gls{bird} backend into the \gls{bra}, the work connects district resources, service requirements, accessibility needs, emissions, and cost within a single tradeoff-oriented optimization pipeline. The Framingham case study demonstrates that this framework can translate real district data into feasible alternative routing designs and Pareto fronts that make the consequences of different planning priorities explicit. For planners, this enables a more transparent form of analysis: instead of asking only for the lowest-cost routes, districts can evaluate how operational, environmental, and equity objectives interact before selecting routes for further testing or implementation.

\section{Future Work}

Future work should first focus on turning the Framingham case study from a planning proof of concept into an operationally calibrated routing tool. This requires replacing the constant-speed travel-time assumption with travel times derived from \gls{fps} \gls{gps} or \gls{avl} traces, time-dependent routing engines, or both. It also requires a unified student-level dataset that links home location, assigned school, grade, eligibility status, current stop, current bus assignment, monitor requirement, wheelchair-accessibility requirement, and district-specific service constraints under an appropriate privacy agreement. With these data, the model could validate current stop assignments, distinguish between route-level and student-level service changes, and test generated routes with operators before any live deployment.

A second direction is to embed equity more directly into the optimization and co-design layers. The current analysis evaluates subgroup outcomes after routing and shows that estimated mean ride times improve for several vulnerable groups, but it does not yet constrain the optimization to guarantee equitable outcomes. Future versions should incorporate stakeholder-defined equity metrics such as subgroup-specific ride-time ceilings, minimum service coverage, accessibility thresholds, concentration-index penalties, or Rawlsian/leximax objectives. The same framework should also be tested beyond Framingham, both on synthetic benchmark districts and on larger real districts, to evaluate scalability, sensitivity to local policy choices, and the degree to which the co-design representation generalizes from school bus routing to broader urban bus network redesign.